\documentclass[11pt]{article}

% Standalone adaptation of the compact classroom-activity typography.
\usepackage[letterpaper,margin=0.25in]{geometry}
\usepackage[T1]{fontenc}
\usepackage{microtype}
\usepackage{amsmath,amssymb}
\usepackage{enumitem}
\usepackage{needspace}
\usepackage{multicol}
\usepackage{xcolor}

\pagestyle{empty}
\linespread{1.04}
\setlength{\parindent}{0pt}
\setlength{\parskip}{4pt}
\setlength{\columnsep}{0.25in}
\setlength{\emergencystretch}{2em}
\setlist[enumerate]{leftmargin=*,itemsep=5pt,topsep=4pt,parsep=0pt}
\AtBeginDocument{
  \setlength{\abovedisplayskip}{5pt plus 1pt minus 1pt}
  \setlength{\belowdisplayskip}{5pt plus 1pt minus 1pt}
  \setlength{\abovedisplayshortskip}{3pt plus 1pt}
  \setlength{\belowdisplayshortskip}{3pt plus 1pt}}
\newcommand{\stageheading}[1]{%
  \par\Needspace{6\baselineskip}%
  \addvspace{.65\baselineskip}%
  {\centering\large\bfseries #1\par}%
  \nopagebreak[4]}
\newcommand{\solution}[1]{%
  \par\addvspace{4pt}\noindent\begin{minipage}{\linewidth}%
  \setlength{\parskip}{4pt}%
  \addvspace{.4\baselineskip}%
  \noindent{\setlength{\fboxsep}{2pt}%
    \colorbox{black!10}{\bfseries\strut Problem #1}}%
  \par\nopagebreak[4]}
\clubpenalty=10000
\widowpenalty=10000
\displaywidowpenalty=10000

\begin{document}

\vspace*{0pt}
{\centering\Large\bfseries Permutations and Combinations I\par}

\stageheading{1. Progressive Problems}

\begin{enumerate}[label=\textbf{\arabic*.}]
\item Six books have different titles. In how many ways can all six be
placed in a single row on a shelf? Each book is used once, and two rows
are different if their left-to-right orders differ.

\item A club has nine distinct members. In how many ways can it choose
a committee of four different members? There are no assigned roles, and
the order in which members are chosen does not matter.

\item Eight distinct students enter a contest with three awards: first,
second, and third place. In how many ways can the awards be assigned if
each award goes to one student and no student receives more than one award?

\item In how many different left-to-right strings can all the letters of
\texttt{BALLOON} be arranged? Use each letter exactly as often as it appears
in the word. Copies of the same letter are indistinguishable, and the
strings need not be English words.

\item Six distinct books, labelled $A,B,C,D,E,F$, are placed in a row,
each exactly once. How many left-to-right arrangements have $A$ and $B$
next to each other, but $C$ and $D$ not next to each other? Either order
within an adjacent pair is allowed.

\item Four girls, Ada, Bella, Cleo, and Dina, and three boys, Eli, Finn,
and Gus, are available for a photograph. Choose five different students
and arrange them in a row. The row must include Ada, at least two girls,
and at least two boys, and no two boys may stand next to each other.
How many different left-to-right rows are possible? All seven students
are distinguishable, and each selected student appears exactly once.
\end{enumerate}

\stageheading{2. Olympiad Challenges}

\begin{enumerate}[label=\textbf{\arabic*.},start=7]
\item Three different numbers are chosen from $\{1,2,\ldots,10\}$.
No two chosen numbers may differ by $1$. How many unordered sets can be
chosen? Each number can be used at most once.

\item Arrange the numbers $1,2,\ldots,7$ in a row, each exactly once.
Call an index $i\in\{1,2,\ldots,6\}$ a \emph{success} if $i$ is
immediately followed by $i+1$, in that order. How many permutations
have exactly two successes? Successes may overlap: for instance,
the consecutive segment $3,4,5$ contributes successes at both $i=3$
and $i=4$. Two permutations are different if their left-to-right orders differ.
\end{enumerate}

\clearpage
\stageheading{3. Solutions}
\begingroup\small
\begin{multicols}{2}

\solution{1}
Order matters because the shelf has a left-to-right order. There are six
choices for the first position, five for the second, and so on. Every
book is distinct and used once, so there is no division for identical
objects and no repetition of choices. The multiplication principle gives
\[
6\cdot5\cdot4\cdot3\cdot2\cdot1=6!=\boxed{720}.
\]

\end{minipage}

\solution{2}
A committee is an unordered group. First count ordered selections of
four different members: there are $9\cdot8\cdot7\cdot6$ choices.
Every committee appears $4!$ times in this count, once for each order of
its four members. Dividing by this common overcount gives
\[
\binom94=\frac{9\cdot8\cdot7\cdot6}{4!}=\boxed{126}.
\]
The members are distinguishable; the division removes selection order,
not differences between members. No member is selected twice.

\end{minipage}

\solution{3}
The awards are distinct positions, so order matters: exchanging the first-
and second-place winners changes the result. Choose the first-place
winner in eight ways, the second in seven, and the third in six.
The decreasing numbers of choices enforce the rule that a student cannot
win twice. Thus
\[
8\cdot7\cdot6=\frac{8!}{5!}=\boxed{336}.
\]
There is no division by $3!$, since the three award assignments are distinct.

\end{minipage}

\solution{4}
There are seven letter positions. Temporarily distinguish the two L's
as $L_1,L_2$ and the two O's as $O_1,O_2$. These seven labelled objects
have $7!$ ordered arrangements, with no object used twice.

In an actual string, exchanging the labels on the L's or on the O's
makes no visible change. Every actual string therefore corresponds to
exactly $2!\,2!$ labelled arrangements. All other letters occur once, so
\[
\frac{7!}{2!\,2!}=\boxed{1260}.
\]
Order still matters; only exchanges of indistinguishable copies are removed.

\end{minipage}

\solution{5}
Count all rows with $A,B$ together, then subtract those in which $C,D$
are also together. Treat $A,B$ as one block. Together with $C,D,E,F$,
this gives five distinct objects, arranged in $5!$ ways. The block can
be $AB$ or $BA$, giving $2\cdot5!$ rows.

For the unwanted rows, treat both pairs as blocks. The four objects
are the two blocks, $E$, and $F$. Each block has two internal orders,
so there are $2\cdot2\cdot4!$ unwanted rows. Hence the required count is
\[
2\cdot5!-4\cdot4!=240-96=\boxed{144}.
\]
Each row has a unique decomposition into these blocks, even when the
blocks touch, so there is no overcount. All six books remain distinct
and appear once; both their order and each block's internal order matter.

\end{minipage}

\solution{6}
Since five students are selected with at least two of each gender, there
are exactly two possible compositions: two girls and three boys, or
three girls and two boys. These cases are disjoint and exhaust the possibilities.
In each case, choose the students first, then arrange them; no student
is chosen twice.

\textbf{Case 1: Two girls and three boys.}
Ada must be one of the girls, so choose her companion in $3$ ways.
All three boys must be chosen. Arrange the two girls in $2!$ ways.
They determine three gaps: before, between, and after them. To separate
three boys, put exactly one in every gap, in $3!$ orders. This gives
\[
3\cdot2!\cdot3!=36.
\]

\textbf{Case 2: Three girls and two boys.}
Choose two of the other three girls to join Ada, and choose two of the
three boys. Arrange the girls in $3!$ ways. Their row has four gaps;
choose two different gaps for the boys and assign the boys to them
in $2!$ orders. This gives
\[
\binom32\binom32\,3!\binom42\,2!=648.
\]
Using different gaps is exactly what prevents adjacent boys. Conversely,
removing the boys from any valid row uniquely recovers the girls' order,
the occupied gaps, and the boys' order, so every valid row is counted once.
Selection is unordered, but the final positions are ordered and all
students are distinguishable. Adding the cases gives
\[
36+648=\boxed{684}.
\]

\end{minipage}

\solution{7}
\textbf{Key insight: remove the compulsory spaces between chosen numbers.}
Write the three chosen numbers in increasing order as $a<b<c$.
This is a unique way to represent an unordered set; it does not add
an ordering choice. The restriction says $b\geq a+2$ and $c\geq b+2$.

Define
\[
x=a,\qquad y=b-1,\qquad z=c-2.
\]
Then $1\leq x<y<z\leq8$. Thus every allowed set produces an ordinary
three-element subset of $\{1,2,\ldots,8\}$.

This correspondence is reversible. From any $1\leq x<y<z\leq8$,
set $a=x$, $b=y+1$, and $c=z+2$. These numbers lie between $1$ and $10$,
and their successive differences are at least $2$. The two constructions
undo each other, so this is a bijection: nothing is lost or counted twice.
All numbers are distinct and none is repeated. Therefore the answer is
\[
\binom83=\frac{8\cdot7\cdot6}{3!}=\boxed{56}.
\]

\end{minipage}

\solution{8}
\textbf{Key insight: required consecutive links can be contracted into blocks,
even when they overlap.}
For each $i\in\{1,\ldots,6\}$, let $E_i$ mean that $i$ is immediately
followed by $i+1$. If a specified set of $k$ links is required, merge
each chain of those links into one block. For example, requiring $E_3$
and $E_4$ creates the single block $(3,4,5)$.

The links follow the path $1,2,\ldots,7$, so each required link reduces
the number of blocks by one. There are exactly $7-k$ distinct blocks,
each with a fixed internal order. Their $(7-k)!$ permutations give all
rows satisfying the required links. Extra successes are allowed in this
intermediate count. Every number is used once, so no indistinguishability
factor is needed.

\textbf{Require two specified successes and exclude the other four.}
Fix a two-element set $S$ of success indices. There are $5!$ permutations
satisfying its two required links. To make these the \emph{only} successes,
apply inclusion--exclusion to the four indices outside $S$:
\[
\begin{gathered}
5!-\binom41 4!+\binom42 3!-\binom43 2!+\binom44 1!\\
=120-96+36-8+1=53.
\end{gathered}
\]
Indeed, requiring any additional $j$ of the four links leaves
$(5-j)!$ block permutations, and there are $\binom4j$ ways to specify them.
The contraction argument also covers overlapping links, so these counts
apply to every choice of $S$.

To see why the correction works, a row with $r$ additional successes
has total weight
\[
\sum_{j=0}^{r}(-1)^j\binom rj=(1-1)^r.
\]
This is $0$ if $r>0$, while a row with no additional successes has weight
$1$. Thus exactly the desired rows survive.

There are $\binom62$ choices for $S$. Every row with exactly two successes
has one unique such set, so multiplying introduces no overcount. The answer is
\[
\binom62\cdot53=15\cdot53=\boxed{795}.
\]

\end{minipage}
\end{multicols}
\endgroup
\end{document}
